\par\medskip\noindent
3.4.8. Consideremos ahora una sucesión exacta de complejos de cadenas
\begin{equation}\tag{3.4.8.1}
 0\longrightarrow L_\bullet\xrightarrow{f}L'_\bullet
 \xrightarrow{g}L''_\bullet\longrightarrow0.
\end{equation}
Esta da funtores de la forma (3.4.3),
\[
 \underline\Psi_{L''_\bullet}\xrightarrow{g^*}
 \underline\Psi_{L'_\bullet}\xrightarrow{f^*}
 \underline\Psi_{L_\bullet}.
\]
La propiedad de transitividad implícita en los funtores de la forma
(3.4.3) da un isomorfismo de funtores
\[
 f^*g^*\simeq(gf)^*=(0_{L_\bullet\to L''_\bullet})^*.
\]

%%%%%%%%%%  scan idx 206  |  folio 195  %%%%%%%%%%

El último miembro es evidentemente isomorfo de manera canónica al funtor
cero cocartesiano de $\underline\Psi_{L''_\bullet}$ a
$\underline\Psi_{L_\bullet}$. Por tanto, el isomorfismo precedente
define un funtor canónico
\begin{equation}\tag{3.4.8.2}
 \underline\Psi_{L''_\bullet}\longrightarrow
 \underline\Psi_{[L_\bullet\longrightarrow L'_\bullet]}.
\end{equation}
Aquí $\underline\Psi_{[L_\bullet\longrightarrow L'_\bullet]}$ denota
la categoría de las parejas $(X',t)$ en las que $X'$ es un objeto de
$\underline\Psi_{L'_\bullet}$ y $t$ es una trivialización de su imagen
inversa $f^*(X')$, es decir, un isomorfismo
\[
 f^*(X')\simeq\sisi{0_{L_\bullet,A'}}{0_{L_\bullet,A}},
\]
donde $0_{L_\bullet,A}$ denota el objeto cero de
$\underline\Psi_{L_\bullet}(A)$ y $A$ es el objeto de
$\underline C$ subyacente a $X'$. Afirmamos que \underline{el
funtor} (3.4.8.2) \underline{es una equivalencia de categorías}.
Junto con la exactitud por la izquierda en $A$ establecida en 3.3.1,
esto completa la propiedad de exactitud en $L_\bullet$ de la expresión
$\underline\Psi_{L_\bullet}(A)$.

Esbocemos la demostración. Dar un objeto del destino de (3.4.8.2)
equivale a dar una extensión
\[
 0\longrightarrow A\xrightarrow{i'}X'\xrightarrow{j'}L'_0
 \longrightarrow0
\]
de $L'_0$ por un objeto $A$ de $\underline C$, junto con un morfismo
$\alpha':L'_1\longrightarrow X'$ que satisface
$\sisi{j}{j'}\alpha'=d'_0$ y $\alpha'd'_1=0$, y un morfismo
$t:L_0\longrightarrow X'$ que satisface $\sisi{j}{j'}t=f_0$ y
$td_0=\alpha'f_1$. Estos datos se representan en el diagrama siguiente:
\begin{center}
\begin{tikzpicture}[
  x=1cm,y=1cm,>=Latex,
  every node/.style={font=\small,inner sep=1.2pt}
]
  \node (z00) at (0,0) {$0$};
  \node (l0)  at (1.3,0) {$L_0$};
  \node (lp0) at (2.8,0) {$L'_0$};
  \node (lpp0) at (4.3,0) {$L''_0$};
  \node (z05) at (5.5,0) {$0$};

  \node (z10) at (0,-1.35) {$0$};
  \node (l1)  at (1.3,-1.35) {$L_1$};
  \node (lp1) at (2.8,-1.35) {$L'_1$};
  \node (lpp1) at (4.3,-1.35) {$L''_1$};
  \node (z15) at (5.5,-1.35) {$0$};

  \node (z20) at (0,-2.70) {$0$};
  \node (l2)  at (1.3,-2.70) {$L_2$};
  \node (lp2) at (2.8,-2.70) {$L'_2$};
  \node (lpp2) at (4.3,-2.70) {$L''_2$};
  \node (z25) at (5.5,-2.70) {$0$};

  \node (xp) at (3.25,1.25) {$X'$};
  \node (ap) at (4.0,2.25) {$A$};
  \node (zap) at (4.75,3.25) {$0$};
  \node (xpp) at (5.85,1.25) {$X''$};
  \node (app) at (6.60,2.25) {$A$};
  \node (zapp) at (7.35,3.25) {$0$};

  \draw[->] (z00)--(l0);
  \draw[->] (l0)--node[above]{$f_0$}(lp0);
  \draw[->] (lp0)--node[above]{$g_0$}(lpp0);
  \draw[->] (lpp0)--(z05);
  \draw[->] (z10)--(l1);
  \draw[->] (l1)--node[above]{$f_1$}(lp1);
  \draw[->] (lp1)--node[above]{$g_1$}(lpp1);
  \draw[->] (lpp1)--(z15);
  \draw[->] (z20)--(l2);
  \draw[->] (l2)--(lp2);
  \draw[->] (lp2)--(lpp2);
  \draw[->] (lpp2)--(z25);

  \draw[->] (l1)--node[right]{$d_0$}(l0);
  \draw[->] (lp1)--node[right]{$d'_0$}(lp0);
  \draw[->] (lpp1)--node[right]{$d''_0$}(lpp0);
  \draw[->] (l2)--(l1);
  \draw[->] (lp2)--node[right]{$d'_1$}(lp1);
  \draw[->] (lpp2)--node[right]{$d''_1$}(lpp1);

  \draw[->] (zap)--(ap);
  \draw[->] (ap)--node[left]{$i'$}(xp);
  \draw[->,dashed] (zapp)--(app);
  \draw[->,dashed] (app)--node[left]{$i''$}(xpp);
  \draw[->] (l0) to[bend left=22] node[above]{$t$} (xp);
  \sisi{\draw[<->] (xp)--node[left]{$j'$}(lp0);}{\draw[->] (xp)--node[left]{$j'$}(lp0);}
  \draw[->] (lp1)--node[pos=.82,right]{$\alpha'$}(xp);
  \draw[->,dashed] (xpp)--node[above]{$j''$}(lpp0);
  \draw[->,dashed] (lpp1) to[bend right=25]
    node[pos=.82,right]{$\alpha''$} (xpp);

  \node[anchor=west] at (7.55,-0.70) {$
    \left\{
    \begin{aligned}
      j'\alpha'&=d'_0,\\
      \alpha'd'_1&=0,\\
      j't&=f_0,\\
      td_0&=\alpha'f_1.
    \end{aligned}
    \right.$};
\end{tikzpicture}
\end{center}

%%%%%%%%%%  scan idx 207  |  folio 196  %%%%%%%%%%

La conocida exactitud por la izquierda, en el primer argumento, de las
extensiones de un objeto $L$ por un objeto $A$, aplicada a la primera fila
del diagrama, muestra que un morfismo $t$ que satisface
$\sisi{j}{j'}t=f_0$ permite interpretar $X'$ como la imagen inversa de
una extensión $X''$ de $L''_0$ por $A$. Bajo esta interpretación, $t$
corresponde a la trivialización natural de la imagen inversa de esa
extensión por $g_0f_0=0$.

El morfismo $\alpha'$ que satisface $\sisi{j}{j'}\alpha'=d'_0$ y
$\alpha'f_1=td_0$ se interpreta mediante esencialmente el mismo
diccionario, aplicado ahora a la segunda fila del diagrama, como proveniente
de un morfismo $\alpha'':L''_1\longrightarrow X''$ que satisface
$j''\alpha''=d''_0$. Aquí $\alpha'$ y $\alpha''$ se consideran como
trivializaciones de las imágenes inversas de $X'$ y $X''$ por $d'_0$ y
$d''_0$, respectivamente. Finalmente, la relación $\alpha'd'_1=0$ es
entonces equivalente a $\alpha''d''_1=0$. Así, los datos
$(X',\alpha',t)$, que forman un objeto de
$\underline\Psi_{[L_\bullet\longrightarrow L'_\bullet]}$, equivalen a
los datos $(X'',\alpha'')$, que forman un objeto de
$\sisi{\underline\Psi_{L_\bullet}}{\underline\Psi_{L''_\bullet}}$.
Esto demuestra, en esencia, que (3.4.8.2) es efectivamente una equivalencia
de categorías.

3.4.9. Sigamos considerando la sucesión exacta (3.4.8.1), así como la
sucesión exacta de grupos $\Ext^i$ que define:
\begin{multline}\tag{3.4.9.1}
 0\longrightarrow\Ext^0(L''_\bullet,A)
 \xrightarrow{g^0}\sisi{\Ext^0(L'_\bullet A)}{\Ext^0(L'_\bullet,A)}
 \xrightarrow{f^0}\Ext^0(L_\bullet,A)
 \xrightarrow{\partial^0}\Ext^1(L''_\bullet,A)\\
 \xrightarrow{g^1}\Ext^1(L'_\bullet,A)
 \xrightarrow{f^1}\Ext^1(L_\bullet,A)\longrightarrow\cdots.
\end{multline}
Los isomorfismos canónicos de las formas (3.1.8) y (3.2.4) dan una
interpretación geométrica de los seis primeros términos de esta sucesión
exacta en función de las categorías $\underline\Psi_{L_\bullet}$,
$\underline\Psi_{L'_\bullet}$ y $\underline\Psi_{L''_\bullet}$. La
compatibilidad (3.4.6) proporciona asimismo una interpretación geométrica
de las flechas

%%%%%%%%%%  scan idx 208  |  folio 197  %%%%%%%%%%

$f^0,g^0,f^1,g^1$. Daremos ahora una interpretación análoga de
$\partial^0$, usando el resultado recién demostrado de que (3.4.8.2) es
una equivalencia de categorías.

Sea
\[
 u\in\Ext^0(L_\bullet,A),
 \qquad\text{es decir,}\qquad
 u:L_0\longrightarrow A,\quad ud_0=0.
\]
Consideremos el objeto $(0,\overline u)$ de
$\underline\Psi_{[L_\bullet\to L'_\bullet]}$. Aquí $0$ es el objeto
cero de $\underline\Psi_{L'_\bullet}$, cuya imagen inversa por $f_0$ se
identifica, por tanto, con el objeto cero de
$\underline\Psi_{L_\bullet}$, y $\overline u$ es el automorfismo de este
último definido por $u$ mediante el isomorfismo canónico (3.1.8). Puesto
que (3.4.8.2) es una equivalencia, $(0,\overline u)$ proviene de un objeto
de $\underline\Psi_{L''_\bullet}(A)$, determinado salvo isomorfismo único.
Su clase de isomorfismo es un elemento de
\[
 \Psi^1_{L''_\bullet}(A)\simeq\Ext^1(L''_\bullet,A).
\]
\underline{Afirmamos que este elemento es precisamente} $\partial^0(u)$.

Demostraremos esta compatibilidad reduciéndola a la análoga de
[4, Cor.~5.5], donde debe corregirse el signo como se indicó en 3.3.6.3.
El argumento proporciona también otra demostración de la equivalencia
(3.4.8.2), reduciéndola a [4, Prop.~2.7]. Considérese la categoría opuesta
$\underline C^{\mathrm o}$ e interprétense la construcción de
$\underline\Psi_{L_\bullet}(A)$, para $A$ fijo, y las construcciones de
3.4.8 y de la presente sección en términos de
$\underline C^{\mathrm o}$ y de la categoría $\underline E$ cofibrada
sobre $\underline C^{\mathrm o}$ cuyos objetos son las extensiones del
objeto fijo $A$ de $\underline C^{\mathrm o}$ por un objeto variable $L$
de $\underline C^{\mathrm o}$. Esta es una categoría cofibrada aditiva
sobre $\underline C^{\mathrm o}$, e incluso exacta por la izquierda en
ella (en el sentido de [4, 2.2]).

Un complejo de cadenas $L_\bullet$ de $\underline C$ puede interpretarse
como

%%%%%%%%%%  scan idx 209  |  folio 198  %%%%%%%%%%

un complejo de cocadenas de $\underline C^{\mathrm o}$. Bajo esta
interpretación, la construcción de $\underline\Psi_{L_\bullet}(A)$
coincide con la opuesta de la categoría
$\underline{\widehat E}(L_\bullet)$ en el sentido de [4, 1.8]. El funtor
(3.4.8.2) se identifica entonces con el opuesto del funtor de la forma
[4, (2.1.2)], reemplazando $E$ por su extensión
$\underline{\widehat E}$ a la categoría de los complejos de cocadenas.
Por consiguiente, la afirmación de que este funtor es una equivalencia es
el caso particular [4, 2.7, a)$\Rightarrow$d)].

Por otra parte, la sucesión exacta (3.4.9.1) se identifica con la sucesión
exacta de los grupos $\Ext^i(A,L''_\bullet)$,
$\Ext^i(A,L'_\bullet)$ y $\Ext^i(A,L_\bullet)$, salvo un factor
$(-1)^i$ en $\partial^i$ (Verdier dixit). Análogamente, los isomorfismos
(3.1.8) y (3.2.4) se identifican con los isomorfismos análogos (3.11.5) y
(3.11.4). La compatibilidad afirmada es, por tanto, el caso particular de
[4, Cor.~5.5] corregido como se indicó antes.

\underline{Observación} 3.4.9.2. El redactor, que teme mucho a los signos,
confiesa no tener plena confianza en la demostración precedente: teme que
se hayan omitido signos al pasar de $C$ a $C^{\mathrm o}$. Recomienda al
lector valiente que compruebe directamente la compatibilidad afirmada (o,
respectivamente, su opuesta). Afortunadamente, estas cuestiones de signos
no desempeñarán ningún papel en las exposiciones siguientes.
